\section{실험 설계}
\label{sec:experiments}

본 절은 여덟 개의 독립 실험 나열 대신, 증거 수집기--계약 컴파일러--상태 관찰자--결과 라우터의 네 구성요소가 어떤 입력을 고정하고 어떤 정책을 비교하는지를 세 증거 블록으로 정리한다. 모든 블록에서 CoC 텍스트, 품질 상태, 자차 움직임, 메타데이터는 관측 또는 도출 입력으로 보존하고, 기한ㆍ반복 정책ㆍ응답 탐지 임계값은 연구자가 명시한 감사 의미론으로 다룬다. 따라서 출력의 PASS/FAIL은 선택한 후보 시간 계약의 통과 또는 검토 경로이며, 차량의 물리 안전/위험 레이블이나 폐루프 제어 성능을 뜻하지 않는다.

\subsection{실험 블록 A: 변환·검사기 타당성}
\label{sec:exp-block-a}

블록 A는 \texttt{scene-0001}과 그 결정적 사건열을 고정한 채, 증거 수집기가 보존한 CoC 시점ㆍ품질ㆍ자차 속도와 계약 컴파일러가 만든 \preserve, $D=3$초 계약을 상태 관찰자가 올바르게 소비하는지 확인한다. 기준 입력에는 최초 2.5초 \textsc{Decelerate} 사건, 3.5초와 4.8초의 반복 사건, 4.6597초에 도출한 자차 속도 응답 시작, 낮은 품질 사건이 포함된다. 정상 모델의 의무 생성ㆍ응답ㆍ기한ㆍ반복 시계ㆍ출처ㆍ교착 질의와 함께, 응답 또는 출처 계보를 의도적으로 훼손한 \texttt{deadline\_2s}, \texttt{clock\_reset}, \texttt{no\_response}, \texttt{untrusted\_overwrite} 합성 변이를 비교한다. 출력은 정상 및 변이별 질의 판정 벡터와 알려진 변이에 대한 기대 오라클 일치이며, 결과 라우터는 이를 검사기 내부 타당성 결과로만 분리한다.

같은 블록에서 \texttt{scene-0001-chained}의 \textsc{Follow} $\rightarrow$ \textsc{Decelerate} $\rightarrow$ \textsc{Stop/Hold} $\rightarrow$ \textsc{Recover} 프로토콜도 비교한다. 관측 CoC 감속 요청과 자차 움직임 응답은 고정하지만, \textsc{Stop/Hold/Recover} 시점은 합성 프로토콜 가정으로 명시하고 \texttt{late\_stop}, \texttt{skip\_hold}, \texttt{premature\_recovery}, \texttt{no\_recovery}의 단계ㆍ순서ㆍ복구 결함을 주입한다. 이 비교의 출력은 각 합성 결함에 대응하는 단계, 순서, 기한, 위험 해소 및 교착 질의의 판정이다.

블록 A의 증거는 RQ1의 \texttt{scene-0001} 변환ㆍ알려진 변이 구분과 RQ4의 정형 관찰자 부가가치에만 대응한다. 이는 실제 오류 검출률, 합성 다단계 궤적의 학습용 적합성 또는 차량 안전을 입증하지 않는다. 메타데이터 미러의 13개 체인과 스크립트 일치 집계는 이 블록의 공유 계약 검증에서 제외하고 블록 C의 구현별 기본 동작 기록으로 분리한다.

\subsection{실험 블록 B: 일괄 감사와 민감도}
\label{sec:exp-block-b}

블록 B는 청크 0의 기록된 94개 장면과 403개 CoC 주석, 해당 자차 움직임 궤적을 고정한다. 증거 수집기는 이 입력에서 290개 \texttt{quality\_pass} 주석을 보존한다. 영어 정규식/복합 의도 규칙에서 방향이 명시된 134개 주석 사건별 후보 계약을 포함하고, 모호하거나 비방향적이거나 복합 의도인 156개는 제외한다($290=134+156$). 134건 중 130개는 속도 변화 탐지기 대상이고, 나머지 4개는 사건 시점에 정지 또는 양보 조건을 이미 만족한 \texttt{PASS\_ALREADY\_SATISFIED} 사건이다. 반복 사건은 같은 응답 시작을 공유할 수 있고 상태형 \preserve/\reset 검사는 선택 에피소드에만 적용한다. 기준 $D=3$초와 \preserve 정책을 기준점으로 두고, \preserve/\reset/\dedup 반복 정책, $D=1,2,3,4,5$초 후보 기한, 그리고 창ㆍ속도 변화 임계값ㆍ방향 비율을 각각 3단계로 조합한 27개 탐지기 설정을 비교한다.

일괄 출력은 94개 장면의 감사 라우팅, 기본 설정의 다섯 검토 후보(\texttt{scene-0019}, \texttt{scene-0028}, \texttt{scene-0042}, \texttt{scene-0065}, \texttt{scene-0074}), 반복 의미론 및 후보 기한에 따른 판정 변화, 그리고 130개 탐지 대상에 대한 27설정 민감도이다. 최초 기본 탐지기 후보는 8건이었고, AI 선별 뒤 $v\leq0.5$~m/s인 STOP/YIELD를 이미 충족으로 처리하는 코호트 보정 탐색 규칙을 사후 추가해 \texttt{scene-0046@10000000}, \texttt{scene-0159@4500000}, \texttt{scene-0159@4900000}을 재분류하여 5건이 되었다. 동시에 94개 장면 레이블의 위반 불리언을 기준ㆍ출처 변이ㆍ응답 변이 모델에 직접 직렬화하여 282개 모델과 1,128개 질의의 기대 판정을 확인한다. 기대 판정도 같은 레이블에서 도출되므로 이는 관찰자 수준 시간 변이 검증이 아니라 플래그/질의 직렬화 회귀이다. 관찰자 수준 시간 변이 근거는 \texttt{scene-0001} 모델에 한정한다.

\begin{table}[t]
\centering
\caption{일괄 감사 블록의 기록된 입력과 생성 범위.}
\label{tab:batch-size}
\begin{tabular}{lr}
\toprule
Item & Value \\
\midrule
Scenes & 94 \\
CoC annotations & 403 \\
Quality-passing annotations & 290 \\
Applicable longitudinal contracts & 134 \\
Speed-change detector subset & 130 \\
Serialized cohort variants & 282 \\
Flag/query regression evaluations & 1,128 \\
\bottomrule
\end{tabular}
\end{table}

따라서 블록 B의 출력은 물리적 결과가 아니라 동일 기록에 대해 어떤 감사 정책과 응답 증거 설정이 검토 우선순위를 바꾸는지의 라우팅ㆍ민감도 결과이다. 403개 주석 전체의 안전성, 다섯 후보의 실제 결함, 단일 탐지기 설정의 보편적 타당성, 또는 $D=3$초의 차량 안전 임계값으로 외삽해서는 안 된다. 이 블록은 RQ2에만 대응한다.

\subsection{실험 블록 C: 출처와 에피소드 확장}
\label{sec:exp-block-c}

블록 C는 ``연결할 수 있는가''라는 출처 적격성 판단과 ``연결 후 시간 계약이 통과하는가''라는 사후 감사 판단을 분리한다. 먼저 로컬 CoC-Nusc의 장면별 시간ㆍ자세 초기화와 장면 간 로그/이전ㆍ다음 관계 부재를 고정하고, 모든 장면을 독립 에피소드로 닫는 정책을 오류 차단 시험으로 사용한다. 합성 두 장면 시간선에서는 장면 A의 감사 수용 요청 주석 사건(0ms), 검증되지 않은 경계(1000ms), 장면 B의 응답 증거(1500ms), 기한(3000ms)을 고정한다. 올바른 \texttt{correct\_discard}와 결함 \texttt{buggy\_carryover}를 비교하여, 상태 관찰자가 경계에서 대기 의무를 \texttt{Discarded}로 닫고 교차 경계 완료와 이월을 금지하는지 검사한다. 출력은 원천 완료, 폐기 도달성, 교차 완료 금지, 이월 금지, 교착 질의의 판정이다.

그 다음 공식 mini는 연속성 스키마를 확인하는 용도로만 사용하고, trainval 후보의 적격성은 로컬 심볼릭 링크로 보관한 nuScenes 호환 메타데이터 미러에서 별도로 감사한다. 같은 로그의 인접성, 1초 이하 간격, 양쪽 장면의 로컬 추론과 자차 움직임 존재를 연속성 적격 조건으로 고정해 16개 엄격 쌍과 13개 체인을 구성했다. 다만 보존된 체인 산출물은 공유 계약과 다른 구현을 사용한다. \texttt{generate\_multiscene\_chain\_model.py}는 느린 의도 토큰을 먼저 찾는 부분문자열 파서라서 ``yield ... then accelerate'' 같은 복합 의도를 감속으로 오분류할 수 있고, 응답 탐지기는 $W=1.0$초, $\delta=0.7$~m/s이며 방향 비율 $\rho$를 사용하지 않는다. 이 사건 배열을 \texttt{LoopSource}/\texttt{Monitor} 두 템플릿에 넣은 13체인 결과와 스크립트 비교는 구현 파이프라인이 끝까지 실행되었는지 보는 스모크 테스트로만 보존한다.

블록 C에서 RQ3의 직접 근거는 로컬 장면 독립성 및 경계 오류 차단 시험이다. 로컬 장면 독립성은 연속 주행의 부정 증거가 아니라 출처 부족 시 의무 이월을 차단하는 정책이다. 13체인 산출물의 수치는 공유 계약의 기한ㆍ응답 또는 장기 에피소드 일반화 근거가 아니며, 공식 trainval 시퀀스 인증이나 객체ㆍ위험 연속성 증명에도 사용할 수 없다.

\begin{table*}[t]
\centering
\caption{실험 블록, 연구 질문, 그리고 해석 경계.}
\label{tab:experiment-rq-map}
\begingroup
\scriptsize
\setlength{\tabcolsep}{2.5pt}
\renewcommand{\arraystretch}{0.90}
\begin{tabularx}{\textwidth}{@{}L{0.14\textwidth}L{0.19\textwidth}L{0.23\textwidth}L{0.18\textwidth}L{0.08\textwidth}Y@{}}
\toprule
실험 블록 & 고정 조건 & 조작·비교 조건 & 출력 & 대응 RQ & 해석 범위 \\
\midrule
블록 A: 변환·검사기 타당성 & \texttt{scene-0001}의 CoCㆍ품질ㆍ자차 응답 사건열 & 응답ㆍ출처 변이, 합성 다단계 결함 & 질의/시간 변이 오라클 판정, 관찰자ㆍ교착 명세 범위 & RQ1, RQ4 & 알려진 합성 결함의 구분과 명세 조합만 지지한다. 실제 오류 검출률, 합성 궤적의 학습 적합성, 차량 안전으로 외삽하지 않는다. \\
\midrule
블록 B: 일괄 감사와 민감도 & 94장면ㆍ403주석의 기록 궤적; 290 품질 통과, 134 방향 명시 사건별 후보, 130 탐지 대상 & 반복 정책, $D=1$--$5$초, 27 탐지기 설정, 플래그 직렬화 & 감사 라우팅, 다섯 검토 후보, 정책ㆍ임계값 민감도, 282모델/1,128질의 회귀 결과 & RQ2 & 방향 명시 부분집합과 정책 의존 우선순위만 지지한다. 직렬화 회귀를 오류 검출력으로 해석하지 않는다. \\
\midrule
블록 C: 출처와 에피소드 확장 & 로컬 장면 독립성; 합성 경계 시간선; 메타데이터 미러와 별도 체인 구현 & \texttt{correct\_discard}/\texttt{buggy\_carryover}; 구현별 체인 파이프라인 & 경계 이월 금지 질의; 13체인 기술 집계 & RQ3 & 경계 오류 차단 시험만 RQ 근거이다. 별도 파서/탐지기의 13체인 수치는 기본 동작 기록이며 정량 추론에서 제외한다. \\
\bottomrule
\end{tabularx}
\endgroup
\end{table*}
